\documentclass[11pt]{article}

%------------------------------------------------
% Packages
%------------------------------------------------
\usepackage[a4paper,margin=2.5cm]{geometry}
\usepackage{amsmath,amssymb,amsthm}
\usepackage{mathtools}
\usepackage{enumitem}
\usepackage{hyperref}

\hypersetup{
    colorlinks=true,
    linkcolor=blue,
    citecolor=blue,
    urlcolor=blue
}

\usepackage[linesnumbered,lined,ruled, noend]{algorithm2e}
\newcommand{\algorithmicrequire}{\textbf{Input:}}
\newcommand{\algorithmicensure}{\textbf{Output:}}
\usepackage{algpseudocode}
\usepackage{enumerate}
\usepackage{graphicx,accents}
\makeatletter
\DeclareRobustCommand{\cev}[1]{%
    \mathpalette\do@cev{#1}%
}
\newcommand{\do@cev}[2]{%
    \fix@cev{#1}{+}%
    \reflectbox{$\m@th#1\vec{\reflectbox{$\fix@cev{#1}{-}\m@th#1#2\fix@cev{#1}{+}$}}$}%
    \fix@cev{#1}{-}%
}
\newcommand{\fix@cev}[2]{%
    \ifx#1\displaystyle
        \mkern#23mu
    \else
        \ifx#1\textstyle
            \mkern#23mu
        \else
        \ifx#1\scriptstyle
            \mkern#22mu
        \else
            \mkern#22mu
        \fi
        \fi
    \fi
}
\makeatother

\DeclareMathOperator{\dom}{{dom\,}}
\DeclareMathOperator{\gcore}{{gcore}}
\DeclareMathOperator{\epi}{Epi}
\DeclareMathOperator{\proj}{P}
\DeclareMathOperator{\dist}{d}
\DeclareMathOperator*{\argmin}{argmin}
\DeclareMathOperator{\prox}{prox}
\DeclareMathOperator{\Id}{Id}
\newcommand{\Tbar}{{\overline{T}}}

\renewcommand{\equiv}{:=}
\newcommand{\ucmin}[2]{\underset{#2}{\mbox{minimize}}~#1}
\newcommand{\map}[3]{#1:\,#2\rightarrow #3\,}
\newcommand{\set}[2]{\left\{#1\,\left|\,#2\right.\right\}}
%\newcommand{\extre}{(-\infty,+\infty]}
\newcommand{\extre}{\overline{\mathbb{R}}}
\newcommand{\xbar}{\bar{x}}
\newcommand{\bbR}{\mathbb{R}}
\newcommand{\bbN}{\mathbb{N}}
\newcommand{\Dcal}{{\mathcal D}}
\newcommand{\paren}[1]{\left(#1\right)}
\newcommand{\Nbb}{\mathbb{N}}
\newcommand{\Mtilde}{{\widetilde{M}}}
\newcommand{\Mbar}{{\overline{M}}}
\newcommand{\diff}{d}


%------------------------------------------------
% Title Information
%------------------------------------------------

\title{\bfseries Quantitative Convergence of Proximal Splitting Iterations in Uniformly Convex Metric Spaces}

\author{
Russell Luke\thanks{Institute for Numerical and Applied Mathematics, Georg-August University of Göttingen,
\texttt{r.luke@math.uni-goettingen.de}}
}

\date{Nov 25-27, 2026}

%------------------------------------------------

\begin{document}

\maketitle

\vspace{-1em}

\noindent
\textbf{Keywords:}
optimization, metric space, proximal mapping, operator splitting

\vspace{0.5em}


\vspace{1em}

%%%%%%%%%%%%%%%%%%%%%%%%%%%%%%%%%%%%%%%%%%%%%%%%%%
\section*{Abstract}
We provide sufficient conditions for quantitative
convergence of the iterates of proximal splitting algorithms for minimizing a sum of functions on a metric space.

The setting for our development is a $p$-uniformly convex metric space, denoted $(G,d)$.
Our central goal is to solve the following optimization problem
 \begin{equation}\label{eq:gen prob}
     \argmin_{x\in G} \sum_{j=1}^m f_j(x).
 \end{equation}
For geodesic metric spaces, in the absence of linear structure {\em implicit tools}
using the {\em proximal mapping} \cite{Jost1995}
% \cite{AriLeuLop14, RuiLopNic15, Izuchukwu19, Bacak13, Bacak14, Bacak23, ChoiJi16, Mayer98, Jost1995, LauLuk21, BLL, OhtaPal15, EspFer09}
are well-developed and provide notions of descent based on  {\em gradient flows} \cite{DeGMarTos80, AmbGigSav2005}.
This is applied to evolution systems (see for instance \cite{GigNob21}) and metric measure spaces (see for instance \cite{OhtaPal15}).
Purely implicit methods based on proximal splitting are possible whenever the proximal mapping can be computed
explicitly \cite{Bacak14, LauLuk21, BLL, LukMir26}.

Let
\begin{equation}
\begin{aligned}
\Tbar &\equiv (\tau_m \prox^p_{f_m,\lambda_m} \oplus (1 - \tau_m)\Id)  \circ \cdots \nonumber \\
  &\quad \circ (\tau_2 \prox^p_{f_2,\lambda_2} \oplus (1 - \tau_2)\Id)
    \circ (\tau_1 \prox^p_{f_1,\lambda_1} \oplus (1 - \tau_1)\Id).
\end{aligned}
\end{equation}
Holding the parameters $\tau_j$ and $\lambda_j$ fixed, we characterize the properties of the mapping $\Tbar$ and
determine the convergence properties of fixed point sequences $(x^k)$ under $\Tbar$:
\[
 x^{k+1}= \overline{T} x^k.
\]

The theory does not assume that the functions $f_j$ have common minima, nor does it require vanishing proximal parameters or step sizes.
Our results are stated for general $p$-uniformly convex
spaces with curvature bounded above, and a corollary specializes the main theorem
to Hadamard spaces, where many assumptions for the more general setting can be dropped.
The theory is demonstrated with computation of Fr\'echet means in the
space of SPD matrices with the affine invariant metric (a Hadamard space) and the
sphere with the usual geodesic metric (a CAT($\kappa$) metric space).

%%%%%%%%%%%%%%%%%%%%%%%%%%%%%%%%%%%%%%%%%%%%%%%%%%
\vspace{1em}

\noindent
\textbf{Acknowledgements.}
This work was supported by the Deutsche Forschungsgemeinschaft (DFG, German Research Foundation) – Projektnummer 566257456 and 541767835

%%%%%%%%%%%%%%%%%%%%%%%%%%%%%%%%%%%%%%%%%%%%%%%%%%
\vspace{1em}

\begin{thebibliography}{9}

\bibitem{AmbGigSav2005}
L.~Ambrosio, N.~Gigli, and G.~Savar{\'e}.
\newblock {\em Gradient Flows in Metric Spaces and in the Space of Probability
  Measures}.
\newblock Basel: Birkh\"auser, 1 edition, 2005.

\bibitem{RuiLopNic15}
D.~{Ariza-Ruiz}, G.~{L{\'o}pez-Acedo}, and A.~Nicolae.
\newblock The asymptotic behavior of the composition of firmly nonexpansive
  mappings.
\newblock {\em J Optim Theory Appl}, 167:409--429, 2015.

\bibitem{Bacak14}
M.~Ba{\v c}{\'a}k.
\newblock Computing medians and means in {{Hadamard}} spaces.
\newblock {\em SIAM J. Optim.}, 24(3):1542--1566, 2014.

\bibitem{BalCarLie94}
K.~Ball, E.~A. Carlen, and E.~H. Lieb.
\newblock Sharp uniform convexity and smoothness inequalities for trace norms.
\newblock {\em Invent. Math.}, 115(3):463--482, 1994.

\bibitem{BLL}
A.~B{\"e}rd{\"e}llima, F.~Lauster, and D.~R. Luke.
\newblock {{$\alpha$}}-firmly nonexpansive operators on metric spaces.
\newblock {\em J. Fixed Point Theory Appl.}, 24, 2022.

\bibitem{BurBurIva}
D.~Burago, {\relax Yu}.~Burago, and S.~Ivanov.
\newblock {\em A Course in Metric Geometry}, volume~33 of {\em Grad. {{Stud}}.
  {{Math}}.}
\newblock Providence, RI: American Mathematical Society (AMS), 2001.

\bibitem{Cheeger99}
J.~Cheeger.
\newblock Differentiability of {{Lipschitz}} functions on metric measure
  spaces.
\newblock {\em Geom. Funct. Anal.}, 9(3):428--517, 1999.

\bibitem{DeGMarTos80}
E.~De~Giorgi, A.~Marino, and M.~Tosques.
\newblock Problemi di evoluzione in spazi metrici e curve di massima pendenza.
\newblock {\em Atti Accad. Naz. Lincei, VIII. Ser., Rend., Cl. Sci. Fis. Mat.
  Nat.}, 68:180--187, 1980.

\bibitem{GigNob21}
Nicola Gigli and Francesco Nobili.
\newblock A differential perspective on gradient flows on
  {{\(\mathbf{\mathsf{CAT}} (\kappa)\)}}-spaces and applications.
\newblock {\em J. Geom. Anal.}, 31(12):11780--11818, 2021.

% \bibitem{GooLewLopNic_cat0cc}
% A.~Goodwin, A.~S. Lewis, G.~{L{\'o}pez-Acedo}, and A.~Nicolae.
% \newblock Convex optimization on {{CAT}}(0) cubical complexes.
% \newblock {\em Adv. Appl. Math.}, 165:20, 2025.
%
% \bibitem{GooLewLopNicFoCM}
% A.~Goodwin, A.~S. Lewis, G.~{L{\'o}pez-Acedo}, and A.~Nicolae.
% \newblock Recognizing {{Weighted Means}} in {{Geodesic Spaces}}.
% \newblock {\em Found Comput Math}, September 2025.
%
% \bibitem{GooLewLopNic26}
% A.~Goodwin, A.~S. Lewis, G.~{L{\'o}pez-Acedo}, and A.~Nicolae.
% \newblock Stochastic and incremental subgradient methods for convex
%   optimization on {{Hadamard}} spaces.
% \newblock {\em Math. Program.}, March 2026.
%
% \bibitem{HeiKos98}
% J.~Heinonen and P.~Koskela.
% \newblock Quasiconformal maps in metric spaces with controlled geometry.
% \newblock {\em Acta Math.}, 181(1):1--61, 1998.

\bibitem{Jost1995}
J.~Jost.
\newblock Convex functionals and generalized harmonic maps into spaces of non
  positive curvature.
\newblock {\em Comment. Math. Helv.}, 70(4):659--673, 1995.

\bibitem{LauLuk21}
F.~Lauster and D.~R. Luke.
\newblock Convergence of proximal splitting algorithms in {{CAT}}({{$\kappa$}})
  spaces and beyond.
\newblock {\em Fixed Point Theory Algorithms Sci Eng}, 2021(13), 2021.

% \bibitem{LewLopNic_SIOPT24}
% A.~S. Lewis, G.~{L{\'o}pez-Acedo}, and A.~Nicolae.
% \newblock Basic convex analysis in metric spaces with bounded curvature.
% \newblock {\em SIAM J. Optim.}, 34(1):366--388, 2024.
%
% \bibitem{Loj63}
% Stanis{\l}aw {\L}ojasiewicz.
% \newblock A topological property of real analytic subsets.
% \newblock Equ. {Derivees} partielles, {Paris} 1962, {Colloques} internat.
%   {Centre} nat. {Rech}. sci. 117, 87-89 (1963)., 1963.
%
\bibitem{LukMir26}
D.~R. Luke and M.~Mirhashemi.
\newblock Convergence of proximal splitting algorithms in uniformly convex
  metric spaces.
\newblock arxiv 2605.03484, 2026.

\bibitem{OhtaPal15}
S.~Ohta and M.~P{\'a}lfia.
\newblock Discrete-time gradient flows and law of large numbers in
  {{Alexandrov}} spaces.
\newblock {\em Calc. Var. Partial Differ. Equ.}, 54(2):1591--1610, 2015.

% \bibitem{Plaut92}
% Conrad Plaut.
% \newblock A metric characterization of manifolds with boundary.
% \newblock {\em Compos. Math.}, 81(3):337--354, 1992.
%
% \bibitem{Polyak63}
% B.~T. Polyak.
% \newblock Gradient methods for the minimisation of functionals.
% \newblock {\em Zh. Vychisl. Mat. Mat. Fiz.}, 3:643--653, 1963.
%
% \bibitem{Rademacher1917}
% H.~Rademacher.
% \newblock {\"Uber partielle und totale Differenzierbarkeit von Funktionen
%   mehrerer Variablen und \"uber die Transformation der Doppelintegrale. I, II.}
% \newblock {\em Math. Ann.}, 79:340--359, 1920.
%
\end{thebibliography}

\end{document}
